\documentclass[11pt,a4paper]{article}
\usepackage[top=1.7cm,bottom=1.5cm,left=1.5cm,right=1.5cm,headsep=8pt,footskip=16pt]{geometry}
\usepackage{enumitem}
\usepackage{amssymb}
\usepackage{array}
\usepackage[table]{xcolor}
\usepackage{fancyhdr}
\usepackage{lastpage}
\usepackage{multicol}
\usepackage{polyglossia}
\setdefaultlanguage[calendar=gregorian, numerals=maghrib]{arabic}
\setotherlanguage{english}
\IfFontExistsTF{Amiri}{%
  \newfontfamily\arabicfont[Script=Arabic, Scale=1.08]{Amiri}
}{%
  \newfontfamily\arabicfont[Script=Arabic, Scale=1.08, Path=../fonts/, Extension=.ttf,
    UprightFont=*-Regular, BoldFont=*-Bold, ItalicFont=*-Italic,
    BoldItalicFont=*-BoldItalic]{Amiri}
}
\IfFontExistsTF{AlBattar}{%
  \newfontfamily\khatt[Script=Arabic, Scale=1.15]{AlBattar}
}{%
  \newfontfamily\khatt[Script=Arabic, Scale=1.15, Path=../fonts/,
    BoldFont=ae_AlBattar.ttf]{ae_AlBattar.ttf}
}
\renewcommand{\headrulewidth}{0.8pt}
\renewcommand{\footrulewidth}{0.8pt}
\fancyhf{}
\rhead{\khatt المدرسة العليا للأساتذة ببوزريعة}
\lhead{\khatt ملحقة بومرداس}
\lfoot{\khatt الأستاذ: د. بن عبد الرزاق يوسف}
\cfoot{\thepage/\pageref{LastPage}}
\pagestyle{fancy}
\parindent 0pt
\setlength{\columnsep}{0.6cm}
\setlength{\columnseprule}{0.4pt}
\setlist[enumerate,1]{label=\textbf{\arabic*.}, leftmargin=*, itemsep=5pt, topsep=2pt}
\begin{document}

\linespread{1.08}\selectfont
\begin{center}
{\LARGE\khatt\textbf{الاختبار السابع عشر في الذكاء الاصطناعي}}\\[4pt]
{\normalsize وضعيات إشكالية ــ أسئلة متعددة الاختيارات ــ الفصول من 1 إلى 6}
\end{center}
\vspace{-2pt}\rule{\textwidth}{0.6pt}\\[2pt]
\textbf{التعليمة:} يتضمن الاختبار خمس وضعيات إشكالية. اقرأ كل وضعية بعناية، ثم اختر التحليل أو التصرف الأنسب
من بين الاختيارات الخمسة؛ لكل وضعية إجابة صحيحة واحدة فقط.
اكتب \textbf{حرف} الإجابة الصحيحة في جدول الإجابات في آخر الاختبار.
لكل وضعية 4 نقاط (المجموع 20 نقطة).
\vspace{4pt}

\noindent
{\large\khatt\textbf{الوضعية 1}}\hfill{\small(الفصل 1)}\par\nopagebreak\vspace{3pt}
\setlength{\fboxsep}{7pt}\colorbox{gray!10}{\parbox{\dimexpr\linewidth-14pt}{يخطط أستاذ لأربعة أنشطة: قراءة نموذجية مسجّلة لقصيدة، ورسم توضيحي لمشهد من قصة، وجدول مقارنة بين المدارس الشعرية، واختبار إلكتروني يُصحَّح آليًا. ويسأل زملاءه عن نوع الأداة التوليدية المناسبة لكل نشاط. وهو لا يملك خبرة سابقة في هذه الأدوات، ويريد أن يتجنب استعمال أداة في غير ما صُممت له. أي توزيع هو الأدق؟}}\par\nopagebreak
\begin{itemize}[leftmargin=1.0cm, labelsep=6pt, itemsep=3pt, topsep=4pt, parsep=0pt]
  \item[\textbf{أ)}] محرك بحث للأنشطة الأربعة، لأن الأدوات التوليدية لا تنتج إلا نصوصًا قصيرة، ولا يمكنها إنتاج صوت أو صورة أو جدول أو أسئلة اختبار بأي شكل من الأشكال.
  \item[\textbf{ب)}] أداة نصية واحدة لكل الأنشطة، لأنها تنتج الصوت والصور والجداول والاختبارات بالجودة نفسها، والتمييز بين أنواع الأدوات لم يعد له معنى في الذكاء الاصطناعي التوليدي، ويكفي أن يحسن صياغة طلبه.
  \item[\textbf{ج)}] أداة صوتية للقراءة، وأداة صور للرسم، وأداة نصية للجدول، أما الاختبار الإلكتروني فلا يمكن لأي أداة توليدية المساعدة فيه، لأن الاختبارات تتطلب برمجة لا تنتجها هذه الأدوات.
  \item[\textbf{د)}] أداة توليد صور للجدول، لأن الجدول شكل بصري، وأداة نصية للرسم التوضيحي، لأن الرسم يحتاج إلى وصف دقيق، وأداة صوتية للاختبار الإلكتروني، لأن التلاميذ يجيبون شفهيًا.
  \item[\textbf{هـ)}] أداة تحويل النص إلى صوت للقراءة، وأداة توليد صور للرسم، وأداة نصية للجدول، وأداة نصية تولّد الأسئلة أو الشيفرة لمنصة اختبارات، مع مراجعة كل مخرج قبل استعماله.
\end{itemize}
\par\vspace{6pt}
\noindent
{\large\khatt\textbf{الوضعية 2}}\hfill{\small(الفصل 2)}\par\nopagebreak\vspace{3pt}
\setlength{\fboxsep}{7pt}\colorbox{gray!10}{\parbox{\dimexpr\linewidth-14pt}{لدى فريق تربوي خمسمئة إنشاء لتلاميذ المتوسط، ويريد اكتشاف أنماط الأخطاء الشائعة دون أن تكون لديه قائمة مسبقة بأصنافها. اقترح عضو تصنيف عينة يدويًا وتدريب نموذج موجَّه، واقترح آخر طريقة تجمع النصوص المتشابهة تلقائيًا. ولا يملك الفريق خبرة كافية في تصنيف الأخطاء مسبقًا، ويريد نتائج قابلة للاستثمار في التدريس. أي تقدير هو الأدق؟}}\par\nopagebreak
\begin{itemize}[leftmargin=1.0cm, labelsep=6pt, itemsep=3pt, topsep=4pt, parsep=0pt]
  \item[\textbf{أ)}] التعلم غير الموجَّه يكفي وحده، لأن المجموعات التي يكوّنها النموذج تحمل أسماء الأخطاء تلقائيًا، فلا حاجة إلى أن يفحصها الخبراء أو يسمّوها أو يتحققوا من تماسكها.
  \item[\textbf{ب)}] الطريقتان متطابقتان تمامًا، والفرق بينهما في الاسم فقط، فيُختار أسرعهما دون اعتبار لوجود أصناف مسبقة أو غيابها، لأن النتائج ستكون واحدة في الحالتين.
  \item[\textbf{ج)}] يُستعمل التعلم الموجَّه بتسمية كل إنشاء «جيد» أو «ضعيف» فقط، لأن هذا التصنيف الثنائي يكشف تلقائيًا كل أنماط الأخطاء الشائعة، دون حاجة إلى طريقة تجمع النصوص المتشابهة.
  \item[\textbf{د)}] لاكتشاف أنماط غير معروفة مسبقًا يناسب التعلم غير الموجَّه الذي يجمع المتشابه دون تسميات، ثم يفحص الخبراء المجموعات ويسمّونها؛ أما الموجَّه فيناسب لاحقًا عندما تتحدد الأصناف.
  \item[\textbf{هـ)}] التعلم الموجَّه هو الوحيد الممكن، لأن النماذج لا تستطيع أن تجمع أي نصوص دون تسميات مسبقة، فيجب أن يخترع الفريق أصنافًا قبل أي تحليل، ولو لم يعرف الأنماط الموجودة فعلًا.
\end{itemize}
\par\vspace{6pt}
\noindent
{\large\khatt\textbf{الوضعية 3}}\hfill{\small(الفصل 3)}\par\nopagebreak\vspace{3pt}
\setlength{\fboxsep}{7pt}\colorbox{gray!10}{\parbox{\dimexpr\linewidth-14pt}{كتب طالب مقالًا نقديًا، وطلب من الأداة «رأيك في مقالي»، فجاء الرد مديحًا عامًا مع ملاحظتين شكليتين. ويريد ملاحظات تساعده فعلًا على تحسين بنائه الحجاجي قبل تسليمه لأستاذه بعد يومين، مع الحفاظ على أسلوبه الشخصي. وقد سمع من زملائه أن الأداة تعيد كتابة النصوص كلها عند طلب التقييم. أي طلب يحقق ذلك؟}}\par\nopagebreak
\begin{itemize}[leftmargin=1.0cm, labelsep=6pt, itemsep=3pt, topsep=4pt, parsep=0pt]
  \item[\textbf{أ)}] «أعد كتابة مقالي ليصبح ممتازًا»، لأن المقال المعاد كتابته يبيّن له كل نقاط الضعف تلقائيًا، ويمكنه تسليم النسخة الجديدة مباشرة لأنها مبنية على أفكاره الأصلية.
  \item[\textbf{ب)}] «كن صريحًا جدًا وقاسيًا في رأيك في مقالي»، لأن القسوة وحدها تكفي لاستخراج ملاحظات عميقة، دون حاجة إلى تحديد معايير التقييم أو منع الأداة من إعادة كتابة المقال، فالنقد اللاذع يكشف الضعف بسرعة.
  \item[\textbf{ج)}] «قيّم مقالي وفق معايير: وضوح الأطروحة، وترابط الحجج، وكفاية الأدلة، ومعالجة الرأي المخالف؛ اذكر لكل معيار موضع ضعف بعينه مع اقتراح، دون إعادة كتابة المقال»، ثم يقرر بنفسه ما يأخذ به.
  \item[\textbf{د)}] «قيّم مقالي وفق معايير محددة واذكر مواضع الضعف»، ثم يقبل كل الملاحظات دون نقاش، لأن الأداة أقدر منه على تقييم البناء الحجاجي، وأي رفض لملاحظاتها يُضعف المقال، ويضمن له علامة أعلى.
  \item[\textbf{هـ)}] «قيّم مقالي وأعطه علامة من عشرين»، لأن العلامة الرقمية أوضح من الملاحظات، ويكفي أن يرفع العلامة بإعادة الطلب حتى يحصل على علامة مرتفعة قبل تسليم المقال.
\end{itemize}
\par\vspace{6pt}
\noindent
{\large\khatt\textbf{الوضعية 4}}\hfill{\small(الفصل 4)}\par\nopagebreak\vspace{3pt}
\setlength{\fboxsep}{7pt}\colorbox{gray!10}{\parbox{\dimexpr\linewidth-14pt}{تحتاج جمعية أولياء التلاميذ ملصقًا لحفل نهاية السنة، فطلب أستاذ من أداة توليد صور رسمًا «بأسلوب الفنان الجزائري فلان» وهو فنان حي معروف، ثم أراد بيع نسخ منه باسم الفنان لتمويل الحفل. ويرى أن الأمر مشروع لأن «الأداة هي التي رسمت». وأيّده بعض الأولياء لأن الهدف خيري. ما التقدير الأدق؟}}\par\nopagebreak
\begin{itemize}[leftmargin=1.0cm, labelsep=6pt, itemsep=3pt, topsep=4pt, parsep=0pt]
  \item[\textbf{أ)}] تقليد أسلوب فنان حي باسمه واستغلاله تجاريًا يثير إشكالًا أخلاقيًا وقد يمس حقوقه؛ فالأسلم وصف الأسلوب دون انتحال اسم الفنان، أو التواصل معه، وعدم الإيحاء بأن العمل من إنتاجه.
  \item[\textbf{ب)}] يكفي أن يكتب في زاوية الملصق «مولَّد بالذكاء الاصطناعي» ليصبح بيعه باسم الفنان مشروعًا، لأن الإفصاح عن الأداة يرفع كل حقوق الفنان على أسلوبه واسمه.
  \item[\textbf{ج)}] يُمنع تمامًا استعمال أي صورة مولَّدة في نشاط مدرسي، حتى لو لم تحاكِ أي فنان، لأن كل صورة مولَّدة تنتهك حقوق فنان ما بالضرورة مهما كان الطلب عامًا، ولو كانت الصورة لأغراض تربوية خالصة.
  \item[\textbf{د)}] الإشكال الوحيد في البيع، أما ذكر اسم الفنان في الملصق المجاني فمقبول دائمًا، لأن الاستعمال غير التجاري يبيح نسبة أي عمل مولَّد إلى أي فنان حي.
  \item[\textbf{هـ)}] لا إشكال إطلاقًا، لأن الأسلوب الفني لا يملكه أحد، والأداة هي التي رسمت، فيمكن بيع الملصق مع ذكر اسم الفنان عليه لزيادة قيمته التجارية دون استئذانه، فالشهرة تخدم الفنان أيضًا.
\end{itemize}
\par\vspace{6pt}
\noindent
{\large\khatt\textbf{الوضعية 5}}\hfill{\small(الفصل 5)}\par\nopagebreak\vspace{3pt}
\setlength{\fboxsep}{7pt}\colorbox{gray!10}{\parbox{\dimexpr\linewidth-14pt}{يبحث طالب عن دراسات حول «أثر القراءة الرقمية في الفهم القرائي»، فطلب من أداة ذكاء اصطناعي ملخصًا للموضوع مع مراجع، فجاءت المراجع مشكوكًا في وجودها. واقترح أستاذه استعمال الأداة بطريقة أخرى في مرحلة البحث. ولا يعرف الطالب كيف يستفيد من الأداة دون الوقوع في المراجع المختلقة مرة أخرى، وموعد التسليم قريب. ما الاستعمال الأنسب؟}}\par\nopagebreak
\begin{itemize}[leftmargin=1.0cm, labelsep=6pt, itemsep=3pt, topsep=4pt, parsep=0pt]
  \item[\textbf{أ)}] يكرر طلب المراجع حتى تتطابق القائمتان، لأن تكرار العناوين نفسها دليل على وجودها، ثم يستشهد بها مباشرة دون قراءة الدراسات، لأن الملخص الذي قدّمته الأداة يكفي.
  \item[\textbf{ب)}] يستعين بالأداة في تحديد الكلمات المفتاحية ومرادفاتها بالعربية والإنجليزية وصياغة استعلامات البحث، ثم يبحث بنفسه في قواعد البيانات العلمية، ويقرأ الدراسات الأصلية قبل الاستشهاد بها.
  \item[\textbf{ج)}] يستعين بالأداة في الكلمات المفتاحية، ثم يطلب منها الحكم على جودة كل دراسة يجدها، ويعتمد حكمها دون قراءة الدراسة، لأن الأداة تقيّم المنهجية أدق من الطالب.
  \item[\textbf{د)}] يطلب من الأداة ملخصات للدراسات بدل عناوينها، ويستشهد بالملخصات في بحثه، لأن ملخص الأداة يغني عن العودة إلى الدراسة الأصلية ما دام الموضوع واضحًا.
  \item[\textbf{هـ)}] يتخلى عن الأداة وعن قواعد البيانات معًا، ويكتفي بالمقالات التي تظهر في الصفحة الأولى لمحرك البحث العام، لأنها الأكثر قراءة، فهي الأكثر موثوقية بالضرورة.
\end{itemize}
\par\vspace{6pt}
\begin{center}
{\khatt\textbf{جدول الإجابات}}\\[3pt]
\renewcommand{\arraystretch}{1.6}
\begin{tabular}{|>{\centering\arraybackslash}m{1.9cm}|*{5}{>{\centering\arraybackslash}m{1.2cm}|}}
\hline
\textbf{رقم السؤال} & \textbf{1} & \textbf{2} & \textbf{3} & \textbf{4} & \textbf{5} \\ \hline
\textbf{الإجابة} &  &  &  &  &  \\[8pt] \hline
\end{tabular}

\vspace{8pt}
\renewcommand{\arraystretch}{2.2}
\begin{tabular}{|>{\centering\arraybackslash}m{7.5cm}|>{\centering\arraybackslash}m{3.2cm}|>{\centering\arraybackslash}m{3.2cm}|}
\hline
\textbf{الاسم واللقب} & \textbf{الفوج} & \textbf{العلامة} \\ \hline
 & & \hfill /20 \\ \hline
\end{tabular}
\end{center}
\end{document}
